Dense-subgraph analytics on graphs is now a routine component of data-management systems, and among the cohesive-subgraph models in production use the $k$-truss occupies a middle ground between the component-oriented $k$-core and the clique: an edge of trussness~$k$ lies in at least $k-2$ triangles whose remaining edges are at least as cohesive.
The model was introduced by Cohen~\cite{cohen2019trusses}, and today the trussness of every edge underlies community detection~\cite{conte2020truly}, truss-based community search~\cite{huang2016attribute,behrouz2022firmtruss}, graph similarity~\cite{zheng2023trussbased}, and temporal cohesive-subgraph analytics~\cite{lotito2020efficient,saryce2025powerful}.
Because the trussness values are recomputed from scratch in $O(\sum_e \tau_e)$ work by bucket peeling, these applications traditionally materialize the decomposition once and reuse it.

The data graphs of the target applications are not static.
Communication, transaction, and collaboration graphs absorb millions of edge updates per day; a decomposition that lags the current state of the graph yields communities that are out of date precisely when they are consulted.
The natural alternative---re-decomposing after every update---costs a full peeling pass per update and is clearly infeasible at update rates, so a line of work has developed \emph{dynamic} $k$-truss maintenance: algorithms that update the maintained trussness values incrementally under edge insertions and deletions.
The state of the art, exemplified by the star-based maintenance of Sun et al.~\cite{sun2023efficient} and the uncertain-graph extension of Sun et al.~\cite{sun2025probabilistic}, achieves remarkable asymptotic bounds per update, but processes the update stream \emph{one edge at a time}: each update performs its own closure computation and its own peeling, and the per-update work does not amortize across a batch of updates that touch the same dense region.

The one-at-a-time discipline is not an implementation convenience; it is baked into the algorithms.
When two updates arrive together---as they do whenever the data source buffers, checkpoints, or replays changes---each update is processed independently, and the second update re-traverses the very region the first update just re-peeled.
On graphs whose dense part is a single large truss-connected component---the empirically dominant case, from email corpora to co-purchase networks---every update's closure covers the entire component, so the total work of a batch is the component size times the batch size rather than the component size.
The parallel batch-dynamic $k$-core machinery of Liu et al.~\cite{liu2022parallel} demonstrates the payoff of batch awareness for the easier core model, where a constant number of rounds suffices; the truss model resists the same treatment because its peeling order couples edge values in a way that the core's simple min-label propagation does not.

This paper closes that gap for the truss model.
We process the stream in batches of arbitrary interleaved insertions and deletions and maintain the \emph{exact} decomposition (every maintained trussness equals the true trussness of the current graph) with an algorithm whose three stages are each parallel: a support-maintenance stage that updates triangle witness counts for the batch's structural changes, a parallel closure stage that computes a sound superset of the changeable edges from the batch's seeds, and a single-pass region peeling stage that recomputes the region's trussness values exactly, in parallel, with boundary edges simulated by \emph{virtual pops}.
The virtual-pop device is the technical heart of the paper: it lets a peeling that sees only a region of the graph reproduce, edge for edge, the trussness values that a full-graph peeling would produce, because every boundary edge is known to be unchanged and therefore contributes exactly one removal event at its maintained bucket.
We prove the resulting procedure correct---the maintained decomposition is exact after every batch---and show that the region computed by our closure contains the true change set, so the peel never assigns a wrong value.

We implemented the algorithm in C++ with a work-stealing thread pool and validated it with a differential-testing harness that compares the maintained decomposition against a static reference implementation after every batch, on randomized graphs, adversarial streams, and real-world datasets.
Experiments over eight public graphs (25K to 69M edges) and streams of 200K updates show that batching collapses the multi-round fixpoint behavior of naive batch designs into one peeling pass, that the maintained decomposition is exact on every workload we tested, and that the sequential per-edge baseline is up to an order of magnitude slower than our batch algorithm on the same workloads while our algorithm approaches the cost of a single static decomposition on graphs with community structure.
To summarize, this paper makes the following contributions:
\begin{itemize}
\item a batch-aware formulation of exact $k$-truss maintenance with a parallel three-stage algorithm: witness maintenance, mate closure, and single-pass region peeling with virtual boundary pops;
\item a correctness argument showing the closure covers the true change set and the region peel reproduces full-graph peeling values exactly, including the exactly-once witness accounting that makes virtual pops sound;
\item an open-source C++ implementation with a differential-testing harness that validates exactness batch by batch;
\item an experimental evaluation over eight public graphs demonstrating that batching and the single-pass peel jointly deliver near-static maintenance costs on community-structured graphs.
\end{itemize}

The rest of the paper proceeds as follows.
\Cref{sec:background} fixes notation and recalls the peeling view of truss decomposition.
\Cref{sec:related} surveys related work.
\Cref{sec:method-region,sec:method-peel} develop the three algorithmic stages and the correctness argument.
\Cref{sec:parallel} describes the parallelization and implementation, and \Cref{sec:experiments} reports the experimental study.
\Cref{sec:conclusion} concludes.
